% =========================================================================
%  Cover Letter v2 — aeh5879
%  Finite Prevention Windows for HIV Post-Exposure Prophylaxis
%  Submission: Science Advances (transferred from Science)
%  A.C. Demidont, Nyx Dynamics LLC
% =========================================================================
\documentclass[11pt]{letter}

\usepackage[T1]{fontenc}
\usepackage[utf8]{inputenc}
\usepackage{mathptmx}
\usepackage[scaled]{helvet}
\usepackage[letterpaper,margin=1in]{geometry}
\usepackage{setspace}
\usepackage[colorlinks=true,linkcolor=black,citecolor=black,urlcolor=blue]{hyperref}
\onehalfspacing

\signature{A.C. Demidont, DO\\
Nyx Dynamics LLC\\
\texttt{ac.demidont@nyxdynamics.com}\\
ORCID: 0000-0002-9216-8569}

\address{A.C. Demidont, DO\\
Nyx Dynamics LLC\\
Philadelphia, PA 19107}

\begin{document}

\begin{letter}{Erin LaFlame\\
Editorial Coordinator\\
\textit{Science Advances}\\
American Association for the Advancement of Science\\
1200 New York Avenue, NW\\
Washington, DC 20005\\
\texttt{elaflame@aaas.org}}

\opening{Dear Ms.\ LaFlame,}

Please find attached the v2 revision of manuscript \textbf{aeh5879}, ``Finite Prevention Windows for HIV Post-Exposure Prophylaxis: Irreversible Proviral Integration Defines Route-Specific Intervention Limits,'' transferred from \textit{Science}. The package includes the revised main manuscript, supplementary materials, and reconciled figures with corrected source files attached separately for the editorial system's figure-rendering pipeline.

\textbf{Substantive corrections.} As disclosed in our prior correspondence, the v1 manuscript's reported $t_{\text{crit}}$ values corresponded to the model's hand-coded \texttt{seeding\_midpoint} parameters rather than to $t_{\text{crit}}(\eta)$ as defined by Theorem~S3.1(iv). The v2 revision replaces the prior phenomenological logistic model with a multiscale within-host implementation (Phase~1 stochastic tau-leaping at low founder population $+$ Phase~2 deterministic ODE handoff $+$ eclipse-phase fixed-delay buffer matching Perelson 1996 Table~2, $\tau \approx 22$~h). Under the canonical operating point ($V_{0} = 10^{3}$ parenteral, $V_{0} = 1$ mucosal, CV $= 0.3$ on $\{\beta, c, \delta, \alpha\}$), the corrected $t_{\text{crit}}$ values at $\eta = 0.05$ are approximately \textbf{34.5~h (parenteral)} and \textbf{60.5~h (mucosal)}, with a route-compression ratio of $\approx 1.75\times$. This compression value is concordant with the empirical 1.5--2$\times$ range spanning Tsai 1998 (intravenous SIV) and Otten 2000 (intravaginal SHIV) NHP challenge data.

\textbf{Population-level efficacy: structural reframe.} The v1 abstract led with the analytical $F_{\text{access}} \times t_{\text{crit}}$ envelope (``below 10\%''). The v2 abstract leads with point estimates from independent cascade-based modeling (companion analysis): under current US policy, the architectural-barrier model produces $P(R_{0}=0) = 0.003\%$ [95\% CI 0.000--0.006\%]; an independent PWID Monte Carlo simulation produces $0.007\%$. Companion outbreak forecasting projects a $92.7\%$ national 10-year outbreak probability under current policy. The analytical $F_{\text{access}} \times t_{\text{crit}}$ envelope ($\sim 11\%$ at the corrected $t_{\text{crit}}$, with sensitivity range 7.5--15.9\%) is consistent with these point estimates as a conservative upper bound. The qualitative claim that PEP cannot serve as a population-level safety net for PWID under current structural conditions is supported by all three independent approaches.

\textbf{Methods and reproducibility refinements.} Beyond the substantive corrections, the v2 revision incorporates standard reproducibility and rigor items: explicit software-environment specification with pinned package versions, deterministic random-seed documentation, sample-size justification with sensitivity analysis, sex-as-a-biological-variable acknowledgment, U=U scoping (the framework is applicable only to transmission-competent exposures from non-virologically-suppressed sources, citing Rodger 2019 PARTNER2 and Bavinton 2018 Opposites Attract), and clarified AI-tool-use disclosure. Independent reproducibility of all numerical claims from committed code at the cited submission SHA \texttt{37e27ea} was verified prior to revision; the disclosure tag \texttt{disclosure-2026-04-30} (commit \texttt{d047d2d}) is preserved as the editorial transparency snapshot.

\textbf{Reconciled numerical-claims summary.} A reconciled summary of all numerical claims with source files and commit references is provided as \texttt{numerical\_claims\_v2.csv} accompanying this submission. Every numerical claim in the v2 manuscript and supplement is traceable to a specific output file at the cited SHA. The submitted code and data are archived at Zenodo (DOI: \texttt{10.5281/zenodo.20044747}); this DOI reflects a fresh deposit at the v2-submission tag and supersedes any prior deposits referenced in earlier correspondence.

\textbf{Figure-rendering issue.} The original combined-PDF rendering pipeline flattened categorical color encoding in several panels. High-resolution figure files are attached separately with the v2 package. v2 Figures~1A and 2A include a vertical shaded band marking the U=U region ($\log_{10}$ VL $< \log_{10}(200)$), per PARTNER2 / modern operational definition; Figure~3 is unchanged because AIDSVu surveillance data is aggregated across mode of HIV acquisition and cannot be reliably subset to U=U status.

We submit this strengthened version in the spirit of full methodological transparency, recognizing that the integrity of self-disclosed corrections depends on the rigor of the corrected manuscript itself. We appreciate your editorial team's careful handling of the submission and the opportunity to proceed to peer review.

Please do not hesitate to contact me with any questions or if any element of the package requires reformatting for the editorial system.

\closing{Sincerely,}

\end{letter}

\end{document}
